% Shared colors and styles of the figures. Every figure reads its data from figures/data/*.csv, written by
% experiments/make_figures.py (comma separated, header row, nan for missing values: pgfplots skips those points).
\definecolor{cA}{HTML}{2a78d6}\definecolor{cB}{HTML}{eb6834}\definecolor{cC}{HTML}{1baf7a}%
\definecolor{cD}{HTML}{4a3aa7}\definecolor{cE}{HTML}{c03b7a}\definecolor{cF}{HTML}{8a6d1f}\definecolor{cGray}{HTML}{8a8984}%
\definecolor{lvA}{HTML}{c6dbf5}\definecolor{lvB}{HTML}{93bbec}\definecolor{lvC}{HTML}{5e98e0}%
\definecolor{lvD}{HTML}{2a78d6}\definecolor{lvE}{HTML}{174a8c}\definecolor{lvF}{HTML}{1baf7a}%
%
% \namedbars{csv file}{ylabel}{title}: lower bounds of the named states at every level, divided by the reference `ref`
% (printed above each state at height `label_y`), with the two upper bounds as markers. Columns: see named_csv in
% make_figures.py. The state positions (`pos`) leave a gap at x = 10 between the PPT and the NPT states.
\providecommand{\namedbar}[3]{% fill, column, offset
  \addplot[ybar, bar width=0.13, bar shift=0pt, draw=none, fill=#1, area legend]
    table[x expr=\thisrow{pos}+#3, y expr=\thisrow{#2}/\thisrow{ref}] {\namedfile};}
\providecommand{\namedbars}[3]{%
  \def\namedfile{#1}%
  \begin{tikzpicture}
    \begin{axis}[width=\textwidth, height=0.29\textwidth, ymin=0, ymax=1.38, xmin=0.3, xmax=14.7,
      ylabel={#2}, ytick={0,0.25,0.5,0.75,1}, xtick=data, xticklabels from table={#1}{state},
      table/col sep=comma, xticklabel style={rotate=35, anchor=north east, font=\scriptsize, inner sep=1pt},
      tick label style={font=\footnotesize}, label style={font=\small}, title={#3}, title style={font=\small,
      yshift=8pt}, clip=false, legend style={font=\scriptsize, at={(0.5,-0.42)}, anchor=north, legend columns=4,
      draw=none, /tikz/every even column/.append style={column sep=6pt}}, legend cell align=left,
      ytick style={draw=none}, xtick style={draw=none}, ymajorgrids, grid style={cGray!25}]
      \addplot[draw=none, forget plot] table[x=pos, y expr=0] {#1};% the ticks (xtick=data)
      \namedbar{lvA}{lb12}{-0.325}\addlegendentry{LB $(1,2)$}
      \namedbar{lvB}{lb13}{-0.195}\addlegendentry{LB $(1,3)$}
      \namedbar{lvC}{lb22}{-0.065}\addlegendentry{LB $(2,2)$}
      \namedbar{lvD}{lb23}{0.065}\addlegendentry{LB $(2,3)$}
      \namedbar{lvE}{lb33}{0.195}\addlegendentry{LB $(3,3)$}
      \namedbar{lvF}{lb24}{0.325}\addlegendentry{LB $(2,4)$}
      \addplot[only marks, mark=diamond*, mark size=2.2pt, cB, mark options={draw=white, line width=0.4pt}]
        table[x expr=\thisrow{pos}-0.14, y expr=\thisrow{ub_extract}/\thisrow{ref}] {#1};
      \addlegendentry{UB: extracted ($k=m=3$)}
      \addplot[only marks, mark=x, mark size=2.4pt, black, thick]
        table[x expr=\thisrow{pos}+0.14, y expr=\thisrow{ub_hjw}/\thisrow{ref}] {#1};
      \addlegendentry{UB: HJW}
      \addplot[only marks, mark=none, nodes near coords={$\pgfplotspointmeta$}, point meta=explicit symbolic,
        nodes near coords style={anchor=south, font=\tiny, inner sep=1pt}, forget plot]
        table[x=pos, y=label_y, meta=ref_label] {#1};
      \draw[cGray] (axis cs:0.3,1) -- (axis cs:14.7,1);
      \draw[cGray, densely dashed] (axis cs:10,0) -- (axis cs:10,1.38);
      \node[font=\scriptsize, anchor=south] at (axis cs:5,1.38) {PPT (bound entangled)};
      \node[font=\scriptsize, anchor=south] at (axis cs:12.5,1.38) {NPT};
    \end{axis}
  \end{tikzpicture}}
%
% \memcurve{column}{color}{mark}{legend}: one curve of memory.csv, broken between the groups of levels of equal degree m
% (pos 1-2: m = 2, pos 4-6: m = 3, pos 8-11: m = 4).
\providecommand{\memcurve}[4]{%
  \addplot[#2, thick, mark=#3, mark size=1.8pt, restrict x to domain=0:2.5] table[x=pos, y=#1] {\memfile};
  \addlegendentry{#4}
  \addplot[#2, thick, mark=#3, mark size=1.8pt, restrict x to domain=3.5:6.5, forget plot] table[x=pos, y=#1] {\memfile};
  \addplot[#2, thick, mark=#3, mark size=1.8pt, restrict x to domain=7.5:11.5, forget plot] table[x=pos, y=#1] {\memfile};}%
\providecommand{\fulllevel}[1]{\fill[cGray!15] (axis cs:#1-0.4,1e3) rectangle (axis cs:#1+0.4,1e16);}%
\begin{figure}
  \centering
  \def\memfile{figures/data/memory.csv}%
  \begin{tikzpicture}
    \begin{axis}[width=\columnwidth, height=0.8\columnwidth, ymode=log, ymin=1e3, ymax=1e16, xmin=0.4, xmax=11.6,
      xtick=data, xticklabels from table={figures/data/memory.csv}{level}, xlabel={level $(k,m)$},
      ytick={1e3,1e6,1e9,1e12,1e15}, yticklabels={1\,kB,1\,MB,1\,GB,1\,TB,1\,PB}, ylabel={memory},
      xticklabel style={rotate=45, anchor=north east, font=\scriptsize, inner sep=1pt},
      tick label style={font=\footnotesize}, label style={font=\small}, xtick style={draw=none},
      ymajorgrids, grid style={cGray!25}, clip=false, table/col sep=comma,
      legend style={font=\scriptsize, at={(0.5,-0.3)}, anchor=north, draw=none, legend columns=2,
      /tikz/every even column/.append style={column sep=4pt}}, legend cell align=left]
      \fulllevel{2}\fulllevel{6}\fulllevel{11}
      \addplot[draw=none, forget plot] table[x=pos, y expr=1e9] {\memfile};% the ticks (xtick=data)
      \draw[black, densely dashed] (axis cs:0.4,1.6e10) -- (axis cs:11.6,1.6e10)
        node[pos=0, anchor=south west, font=\scriptsize, inner sep=1pt] {16\,GB};
      \memcurve{mem_r4a2}{cA}{*}{$r=4$, $r_A=2$}
      \memcurve{mem_r4a3}{cB}{square*}{$r=4$, $r_A=3$}
      \memcurve{mem_r9a3}{cC}{triangle*}{$r=9$, $r_A=3$}
      \memcurve{mem_r4a4}{cD}{diamond*}{$r=4$, $r_A=4$}
      \memcurve{mem_r8a4}{cE}{pentagon*}{$r=8$, $r_A=4$}
      \memcurve{mem_r16a4}{cF}{star}{$r=16$, $r_A=4$}
      \node[font=\scriptsize, anchor=south] at (axis cs:1.5,1e16) {$m=2$};
      \node[font=\scriptsize, anchor=south] at (axis cs:5,1e16) {$m=3$};
      \node[font=\scriptsize, anchor=south] at (axis cs:9.5,1e16) {$m=4$};
    \end{axis}
  \end{tikzpicture}
  \caption{Memory needed by an interior-point solver for the bound at each level $(k,m)$, as a function of
  $r=\operatorname{rank}\rho$ and $r_A=\operatorname{rank}\rho_A$, for the full hierarchy [$k=m$, shaded,
  Eq.~\eqref{eq:SDPcroof1}] and the simplified one [$k<m$, Eq.~\eqref{eq:marginals_SDP}]. After the reductions of
  Appendix~\ref{app:SymRed}, the SDP only depends on $r$, $r_A$ and $(k,m)$; we count its number $n$ of free real
  variables (after eliminating the permutation-invariance equalities) and show the size $8n^2$ bytes of the dense
  Schur complement matrix, which dominates the memory of the solver for large problems. At equal degree $m$, the
  simplified levels with $k\geq2$ are cheaper than the full level only for states of high rank ($r=9$, $r_A=3$ and
  $r=16$, $r_A=4$); for low-rank states ($r=4$, $r_A=3$ and $r=8$, $r_A=4$) the full level $(m,m)$ needs less memory,
  since then $r<(k+1)r_A-k$ and a full copy costs less than a marginal copy. Dashed: the 16\,GB of a laptop.}
  \label{fig:memory}
\end{figure}
